\chapter{Resolution Limits of Amplitude-Based Migration}
\label{ch:resolution}

This chapter establishes the baseline problem that motivates the rest of the
thesis: how closely can two point scatterers be spaced, in a medium with a
known wave velocity, before standard migration imaging can no longer tell
them apart? The experiment uses two subwavelength perfect-electric-conductor
(PEC) cylinders embedded in ice, illuminated by a zero-offset GPR B-scan
simulated in gprMax, and migrated with three independent algorithms. The
results define the amplitude-based resolution floor against which the
phase-based method of \cref{ch:phase-plane} is later compared.

\paragraph{Hypothesis.} Two point scatterers cannot be resolved as separate
amplitude peaks in a migrated GPR image once their separation drops below
roughly half a wavelength, regardless of which migration algorithm
(Kirchhoff, Gazdag, or back-propagation) is used.

\section{Ricker Wavelet and Model Setup}

The simulated antenna emits a Ricker wavelet with centre frequency
$f_c = \SI{15}{\GHz}$ and a time-zero offset $t_0 = \SI{0.0943}{\ns}$
(\cref{fig:res-setup}a). Two PEC cylinders of radius $r = \SI{28}{mm}$ are
buried at a depth of \SI{0.676}{m} in a $4.01 \times 4.01$ m domain
discretised at $\Dx = \SI{1}{mm}$ with absorbing PML boundaries
(\cref{fig:res-setup}b). The cylinder separation is swept geometrically from
$2\lambda$ down to $\tfrac{1}{16}\lambda$ across a family of scenarios
(\cref{fig:res-setup}c), where $\lambda$ is the dominant wavelength of the
Ricker pulse in ice.

\begin{figure}[htbp]
  \centering
  \begin{subfigure}[t]{0.32\textwidth}
    \includegraphics[width=\linewidth]{RES_001_Ricker_Wavelet_f_c__15_GHz_t0__0943_ns.png}
    \caption{Source wavelet}
  \end{subfigure}
  \hfill
  \begin{subfigure}[t]{0.32\textwidth}
    \includegraphics[width=\linewidth]{RES_002_Resolution_Study__Model_Geometry__domain_4010_m_Δx__1_mm_PML.png}
    \caption{Model geometry}
  \end{subfigure}
  \hfill
  \begin{subfigure}[t]{0.32\textwidth}
    \includegraphics[width=\linewidth]{RES_003_Scatterer_Positions__PEC_Cylinders__r__28_mm_depth__0676_m.png}
    \caption{Scatterer positions}
  \end{subfigure}
  \caption{Forward-model setup for the resolution study: (a) the Ricker
    source wavelet ($f_c=\SI{15}{\GHz}$, $t_0=\SI{0.0943}{\ns}$); (b) the
    gprMax domain and grid; (c) the swept separation between the two PEC
    cylinder scatterers, $r=\SI{28}{mm}$, depth $\SI{0.676}{m}$.}
  \label{fig:res-setup}
\end{figure}

\section{Raw B-Scan Generation}

\Cref{fig:res-bscans} shows the simulated zero-offset B-scans before any
processing: the background-only model, the two-scatterer model for each
separation scenario, and the result of subtracting the background trace from
each trace to suppress the direct air/ground wave and isolate the scatterer
hyperbolae.

\begin{figure}[htbp]
  \centering
  \begin{subfigure}[t]{0.48\textwidth}
    \includegraphics[width=\linewidth]{RES_004_GPR_B-Scans__Background_and_Separation_Models.png}
    \caption{Background and separation models}
  \end{subfigure}
  \hfill
  \begin{subfigure}[t]{0.48\textwidth}
    \includegraphics[width=\linewidth]{RES_005_GPR_B-Scans__Background_Subtracted.png}
    \caption{Background-subtracted}
  \end{subfigure}
  \caption{Raw zero-offset B-scans for the resolution study, before (a) and
    after (b) background subtraction.}
  \label{fig:res-bscans}
\end{figure}

\section{Pre-Migration Signal Conditioning}

Before migration, each trace is conditioned with an exponential decay taper
that suppresses late-arriving energy, a cosine end-taper that removes the
hyperbola tails at the edge of the aperture, and a time-zero shift that
aligns the surface reflection to $t=0$. \Cref{fig:res-taper} shows the effect
of this conditioning on a single trace and on the full B-scan for the
$2\lambda$ separation scenario.

\begin{figure}[htbp]
  \centering
  \begin{subfigure}[t]{0.48\textwidth}
    \includegraphics[width=\linewidth]{RES_006_Effect_of_Tapering_and_t0_Shift__2λ_dataset_single_trace.png}
    \caption{Single trace}
  \end{subfigure}
  \hfill
  \begin{subfigure}[t]{0.48\textwidth}
    \includegraphics[width=\linewidth]{RES_007_B-scan_effect_of_tapering_and_t0_shift__2λ_dataset.png}
    \caption{Full B-scan}
  \end{subfigure}
  \caption{Effect of tapering and the $t_0$ shift on the $2\lambda$ dataset:
    (a) a single representative trace; (b) the complete B-scan.}
  \label{fig:res-taper}
\end{figure}

\section{Kirchhoff Migration}

Kirchhoff (delay-and-sum) migration is applied with the PyLops zero-offset
operator under the exploding-reflector assumption $\vmig = v/2$
(\cref{ch:theory}, \cref{ch:methodology}). \Cref{fig:res-kirchhoff} shows the
migrated image for every separation scenario at $f_c=\SI{15}{\GHz}$ with an
aperture of 40 traces, alongside a zoomed view around the true scatterer
depth.

\begin{figure}[htbp]
  \centering
  \begin{subfigure}[t]{0.48\textwidth}
    \includegraphics[width=\linewidth]{RES_008_Kirchhoff_Migration__All_Datasets____f_c15_GHz____aperture40.png}
    \caption{All separation scenarios}
  \end{subfigure}
  \hfill
  \begin{subfigure}[t]{0.48\textwidth}
    \includegraphics[width=\linewidth]{RES_009_Kirchhoff_Migration_zoomed____f_c15_GHz____aperture40.png}
    \caption{Zoomed around scatterer depth}
  \end{subfigure}
  \caption{Kirchhoff migration of the resolution study ($f_c=\SI{15}{\GHz}$,
    aperture $=40$ traces): (a) all separation scenarios; (b) zoomed view.}
  \label{fig:res-kirchhoff}
\end{figure}

\section{Gazdag Phase-Shift Migration}

The same B-scans are migrated with the Gazdag $f$-$k$ phase-shift algorithm
(downward continuation with evanescent-bin filtering, \cref{ch:theory}).
\Cref{fig:res-gazdag} shows the resulting images and a zoomed view.

\begin{figure}[htbp]
  \centering
  \begin{subfigure}[t]{0.48\textwidth}
    \includegraphics[width=\linewidth]{RES_010_Gazdag_Phase-Shift_Migration__All_Datasets____f_c15_GHz____z.png}
    \caption{All separation scenarios}
  \end{subfigure}
  \hfill
  \begin{subfigure}[t]{0.48\textwidth}
    \includegraphics[width=\linewidth]{RES_011_Gazdag_Phase-Shift_Migration_zoomed____f_c15_GHz.png}
    \caption{Zoomed around scatterer depth}
  \end{subfigure}
  \caption{Gazdag phase-shift migration of the resolution study
    ($f_c=\SI{15}{\GHz}$): (a) all separation scenarios; (b) zoomed view.}
  \label{fig:res-gazdag}
\end{figure}

\section{Back-Propagation Migration}

As a third, independent cross-check, each B-scan is time-reversed and
re-injected into the gprMax domain (\cref{ch:theory}); the migrated image is
read off as the field snapshot at the focusing time
$t_{\mathrm{focus}} = T - t_0$. Both the full electric-field magnitude
$\lVert\mathbf{E}\rVert$ and the single $E_z$ component are recorded, each
with a zoomed view (\cref{fig:res-backprop}).

\begin{figure}[htbp]
  \centering
  \begin{subfigure}[t]{0.48\textwidth}
    \includegraphics[width=\linewidth]{RES_012_Back-Propagation_E__All_Datasets____focus_at_1906_ns.png}
    \caption{$\lVert\mathbf{E}\rVert$, all scenarios}
  \end{subfigure}
  \hfill
  \begin{subfigure}[t]{0.48\textwidth}
    \includegraphics[width=\linewidth]{RES_013_Back-Propagation_E_zoomed____focus_at_1906_ns.png}
    \caption{$\lVert\mathbf{E}\rVert$, zoomed}
  \end{subfigure}

  \vspace{0.6em}
  \begin{subfigure}[t]{0.48\textwidth}
    \includegraphics[width=\linewidth]{RES_014_Back-Propagation_Ez__All_Datasets____focus_at_1906_ns.png}
    \caption{$E_z$, all scenarios}
  \end{subfigure}
  \hfill
  \begin{subfigure}[t]{0.48\textwidth}
    \includegraphics[width=\linewidth]{RES_015_Back-Propagation_Ez_zoomed____focus_at_1906_ns.png}
    \caption{$E_z$, zoomed}
  \end{subfigure}
  \caption{Time-reversal back-propagation migration of the resolution study,
    focused at $t=\SI{1906}{\ns}$: field-magnitude image (a,b) and the $E_z$
    component (c,d), each with a zoomed view around the scatterer depth.}
  \label{fig:res-backprop}
\end{figure}

\section{Comparing Migration Methods and the Amplitude Resolution Limit}
\label{sec:res-comparison}

\Cref{fig:res-psf}a overlays the signed migrated amplitude from all
three algorithms at $f_c=\SI{15}{\GHz}$, and \cref{fig:res-psf}b plots the
normalised lateral point-spread function (PSF) extracted at the true
scatterer depth as a function of separation.

\begin{figure}[htbp]
  \centering
  \begin{subfigure}[t]{0.48\textwidth}
    \includegraphics[width=\linewidth]{RES_016_Migration_Comparison__Signed_Amplitude____f_c15_GHz____apert.png}
    \caption{Signed-amplitude comparison}
  \end{subfigure}
  \hfill
  \begin{subfigure}[t]{0.48\textwidth}
    \includegraphics[width=\linewidth]{RES_017_Normalised_Lateral_PSF_at_True_Scatterer_Depth.png}
    \caption{Normalised lateral PSF}
  \end{subfigure}
  \caption{(a) Signed migrated amplitude for Kirchhoff, Gazdag, and
    back-propagation migration overlaid at $f_c=\SI{15}{\GHz}$; (b) the
    normalised lateral point-spread function at the true scatterer depth,
    swept across separations from $2\lambda$ to $\tfrac{1}{16}\lambda$.}
  \label{fig:res-psf}
\end{figure}

All three migration algorithms collapse the two scatterer hyperbolae into
distinguishable amplitude peaks for separations down to roughly half a
wavelength, but the two peaks progressively merge into a single lobe as the
separation shrinks further: the amplitude image alone cannot certify two
scatterers, or a sub-wavelength displacement of one scatterer, below this
floor. \draftnote{state the precise separation at which the methods stop
resolving two distinguishable PSF peaks, read directly off
\cref{fig:res-psf}b, and comment on any difference between the three
algorithms.} This amplitude-based floor is the motivation for the
phase-plane approach developed in \cref{ch:theory} and validated in
\cref{ch:phase-plane}.
